\section{Estimating the Hardware Speed}\label{app:hardware}

This appendix gives the reasoning behind assumption A2 and the advantage factor $A$ of \cref{sec:pricing}.

\paragraph{SHA-256.}
We recall the parts of SHA-256~\cite{FIPS180-4} that determine the latency of a chain step.

\begin{definition}[SHA-256 compression]\label{def:sha}
Words are 32 bits, $+$ is addition modulo $2^{32}$, and $x^{\ggg r}$ and $x^{\gg r}$ are $x$ rotated and shifted right by $r$ bits. Let
\begin{alignat*}{2}
  \mathrm{Ch}(e,f,g) &= (e \wedge f) \oplus (\neg e \wedge g), &\quad
  \mathrm{Maj}(a,b,c) &= (a \wedge b) \oplus (a \wedge c) \oplus (b \wedge c),\\
  \Sigma_0(a) &= a^{\ggg 2} \oplus a^{\ggg 13} \oplus a^{\ggg 22}, &\quad
  \Sigma_1(e) &= e^{\ggg 6} \oplus e^{\ggg 11} \oplus e^{\ggg 25},\\
  \sigma_0(x) &= x^{\ggg 7} \oplus x^{\ggg 18} \oplus x^{\gg 3}, &\quad
  \sigma_1(x) &= x^{\ggg 17} \oplus x^{\ggg 19} \oplus x^{\gg 10}.
\end{alignat*}
The compression function maps a state $V$ of eight words and a block $M = W_0 \cat \cdots \cat W_{15}$ of sixteen words to a new state. It extends the block by $W_t = \sigma_1(W_{t-2}) + W_{t-7} + \sigma_0(W_{t-15}) + W_{t-16}$ for $16 \le t < 64$.

The rounds start from $(a,b,c,d,e,f,g,h) = V$. For $t = 0, \ldots, 63$, with fixed round constants $K_t$, round $t$ computes
\begin{align*}
  T_1 &= h + \Sigma_1(e) + \mathrm{Ch}(e,f,g) + K_t + W_t, \qquad T_2 = \Sigma_0(a) + \mathrm{Maj}(a,b,c),\\
  (a,b,c,d,e,f,g,h) &\gets (T_1 + T_2,\ a,\ b,\ c,\ d + T_1,\ e,\ f,\ g).
\end{align*}
The output is $V + (a,b,c,d,e,f,g,h)$, added word by word.

For a message $x$ of at most 55 bytes, $\SHA(x)$ is the compression of the fixed initial state with the single block $x \cat \mathtt{0x80} \cat 0^{55-|x|} \cat \mathsf{be64}(8|x|)$.
\end{definition}

% One SHA-256 round as a circuit: textbook adders vs carry-save + retiming (Appendix C).
\begin{figure}[t]
  \centering
  \begin{tikzpicture}[
      x=0.75cm, y=1cm, font=\small,
      fn/.style={draw, rounded corners=2pt, minimum width=1.0cm, minimum height=0.55cm, fill=black!3},
      add/.style={draw, circle, inner sep=0pt, minimum size=0.5cm, fill=black!3},
      csa/.style={draw, rounded corners=2pt, align=center, fill=black!3, font=\footnotesize},
      inp/.style={font=\footnotesize, text=black!70},
      arr/.style={-{Stealth[length=4pt]}, line width=0.5pt},
      hot/.style={-{Stealth[length=4.5pt]}, line width=1.2pt, draw=warn},
      hotbox/.style={draw=warn, line width=1.1pt},
      pre/.style={draw, dashed, rounded corners=2pt, align=center, text=black!60, font=\footnotesize},
      capt/.style={font=\footnotesize\bfseries, anchor=west},
    ]
    % ---------------- (a) textbook datapath ----------------
    \node[capt] at (-1.8,2.25) {(a) Carry-propagate adders only};
    \node[fn, hotbox] (S1) at (0,0) {$\Sigma_1$};
    \node[fn] (Ch) at (0,-0.85) {$\mathrm{Ch}$};
    \node[fn] (S0) at (0,-2.0) {$\Sigma_0$};
    \node[fn] (Mj) at (0,-2.85) {$\mathrm{Maj}$};
    \node[inp, anchor=east] (ie) at (-1.0,0) {$e$};
    \node[inp, anchor=east] at (-1.0,-0.85) {$e,f,g$};
    \node[inp, anchor=east] at (-1.0,-2.0) {$a$};
    \node[inp, anchor=east] at (-1.0,-2.85) {$a,b,c$};
    \draw[hot] (ie) -- (S1);
    \draw[arr] (-1.0,-0.85) -- (Ch);
    \draw[arr] (-1.0,-2.0) -- (S0);
    \draw[arr] (-1.0,-2.85) -- (Mj);
    \node[add, hotbox] (A1) at (1.5,0) {$+$};
    \node[add, hotbox] (A2) at (2.6,0) {$+$};
    \node[add, hotbox] (A3) at (3.7,0) {$+$};
    \node[add, hotbox] (A4) at (4.8,0) {$+$};
    \node[add] (A5) at (1.5,-2.42) {$+$};
    \node[add, hotbox] (A6) at (6.3,0) {$+$};
    \node[add] (A7) at (6.3,0.95) {$+$};
    \draw[hot] (S1) -- (A1);
    \draw[arr] (Ch) -| (A1);
    \draw[hot] (A1) -- (A2);
    \draw[hot] (A2) -- (A3);
    \draw[hot] (A3) -- (A4);
    \draw[arr] (2.6,0.75) node[inp, above] {$h$} -- (A2);
    \draw[arr] (3.7,0.75) node[inp, above] {$K_t$} -- (A3);
    \draw[arr] (4.8,0.75) node[inp, above] {$W_t$} -- (A4);
    \draw[hot] (A4) -- node[inp, below, pos=0.45] {$T_1$} (A6);
    \draw[arr] (5.5,0) |- (A7);
    \draw[arr] (6.3,1.7) node[inp, above] {$d$} -- (A7);
    \draw[arr] (S0) -| (A5);
    \draw[arr] (Mj) -| (A5);
    \draw[arr] (A5) -- node[inp, below, pos=0.5] {$T_2$} (6.3,-2.42) -- (A6);
    \draw[hot] (A6) -- (7.2,0) node[anchor=west, text=warn] {$a'$};
    \draw[arr] (A7) -- (7.2,0.95) node[anchor=west] {$e'$};
    \node[font=\footnotesize, text=warn, align=left, anchor=north west] at (2.0,-0.55)
      {five carry-propagate\\additions in series};

    % ---------------- (b) carry-save + retiming ----------------
    \begin{scope}[xshift=7.7cm]
      \node[capt] at (-1.8,2.25) {(b) Carry-save adders and retiming};
      \node[fn, hotbox] (S1b) at (0,0) {$\Sigma_1$};
      \node[fn, hotbox] (Chb) at (0,-0.85) {$\mathrm{Ch}$};
      \node[fn] (S0b) at (0,-2.0) {$\Sigma_0$};
      \node[fn] (Mjb) at (0,-2.85) {$\mathrm{Maj}$};
      \node[inp, anchor=east] (ieb) at (-1.0,0) {$e$};
      \node[inp, anchor=east] (iefg) at (-1.0,-0.85) {$e,f,g$};
      \node[inp, anchor=east] at (-1.0,-2.0) {$a$};
      \node[inp, anchor=east] at (-1.0,-2.85) {$a,b,c$};
      \draw[hot] (ieb) -- (S1b);
      \draw[hot] (iefg) -- (Chb);
      \draw[arr] (-1.0,-2.0) -- (S0b);
      \draw[arr] (-1.0,-2.85) -- (Mjb);
      \node[pre] (u) at (-0.4,0.95) {$u_t = h{+}K_t{+}W_t$,\\[-2pt]one round early};
      \node[csa, hotbox, minimum width=1.25cm, minimum height=4.1cm] (C) at (2.75,-1.2)
        {carry-\\save\\layer};
      \draw[hot] (S1b) -- (S1b -| C.west);
      \draw[hot] (Chb) -- (Chb -| C.west);
      \draw[arr] (S0b) -- (S0b -| C.west);
      \draw[arr] (Mjb) -- (Mjb -| C.west);
      \draw[arr, dashed] (u.east) -| (1.75,0.4) -- ($(C.west)+(0,1.6)$);
      \draw[arr] (2.75,1.35) node[inp, above] {$d$} -- (C.north);
      \node[csa, hotbox, minimum width=1.6cm, minimum height=0.55cm] (P1) at (5.25,0) {32-bit add};
      \node[csa, hotbox, minimum width=1.6cm, minimum height=0.55cm] (P2) at (5.25,-2.42) {32-bit add};
      \draw[hot] ($(C.east)+(0,1.2)$) -- (P1.west);
      \draw[hot] ($(C.east)+(0,-1.22)$) -- (P2.west);
      \draw[hot] (P1) -- (6.6,0) node[anchor=west, text=warn] {$e'$};
      \draw[hot] (P2) -- (6.6,-2.42) node[anchor=west, text=warn] {$a'$};
      \node[font=\footnotesize, text=warn, align=left, anchor=north west] at (3.85,-0.6)
        {one carry-save\\layer and one\\32-bit add per\\new word};
    \end{scope}
  \end{tikzpicture}
  \caption{One SHA-256 round as a circuit, with the path that limits a chain's latency in red. Inputs on the left are the current state words; $a'$ and $e'$ return to the state registers for the next round. (a)~A direct implementation adds the terms with 32-bit carry-propagate adders, which puts several of them in series. (b)~With carry-save adders and retiming~\cite{LCV07}, $u_t = h + K_t + W_t$ is computed during the previous round, since $h$ is the previous $g$. The remaining terms are reduced to two numbers in a carry-save layer, whose depth does not depend on the word length, and each new word needs one carry-propagate addition. The drawing is schematic; the retimed datapath of~\cite{LCV07} differs in detail.}
  \label{fig:sha-round}
\end{figure}


A chain step $F_i$ hashes 52 bytes, so it is one compression. The message words $W_t$ depend only on the step's input and can be computed alongside the rounds, so the latency of a step is that of the 64 rounds.

\paragraph{How SHA-256 is computed in hardware.}
Assumption A2 can be examined against the structure of SHA-256. One SHA-256 compression is 64 rounds, and each round needs the exact output of the previous one. The new state words are sums of up to five 32-bit terms. Some of them, such as the round constant and the message word, can be added ahead of time, but the next round applies rotations and bitwise functions to the exact bits of the result. Every round must therefore finish at least one 32-bit addition with full carry propagation, plus a few levels of logic. Lee, Chan and Verbauwhede~\cite{LCV07} show that this much is also enough: with carry-save adders and retiming, the loop through one round needs one carry-propagate addition, one carry-save layer and the choice function $\mathrm{Ch}$. \Cref{fig:sha-round} shows both datapaths.

Pipelining raises throughput but not the latency of one chain, since each step starts from the previous step's output. Unrolling removes the registers between rounds, so the relevant floor is the one for logic alone. Hence $\ell$ is roughly 64 times the shortest critical path of one round, measured over many unrolled rounds.

\paragraph{The circuit model.}
The model row of \cref{tab:pricing} quantifies this argument. Its unit is the FO4, the delay of an inverter driving four copies of itself. It estimates the round loop of Lee et al.\ at 10--17 FO4: 7--12 for a fast 32-bit adder~\cite{HS03} and 3--5 that we assume for the carry-save layer and $\mathrm{Ch}$. Predicted FO4 delays for 3--5\,nm processes are 4.3--4.7\,ps~\cite{NCTI21}. Then 64 rounds take
\[
  64 \times 10 \times 4.3\,\text{ps} = 2.75\,\text{ns}
  \quad\text{to}\quad
  64 \times 17 \times 4.7\,\text{ps} = 5.11\,\text{ns},
\]
and with 0.1--0.3\,ns for the final addition and startup, 2.85--5.41\,ns. Our allowance of 1 to 1.8 times for wires, cell choice and registers gives 2.9--9.7\,ns, that is, $A \approx 1.3$--$4.5$.

As a cross-check, Intel's mining patent describes a state update 13 logic levels deep~\cite{SSM18}; if each level cost one FO4, which is optimistic, a hash would take 3.6--3.9\,ns. A process shrink by itself gives only $A \approx 1.3$. At 3.5\,ns a round has 55\,ps, about 12 FO4, so the design point needs a highly optimized custom datapath and sits near the aggressive end of the model. The model's lower end is a conditional estimate, not a physical bound, and no published design that we know of comes close to it.

\paragraph{Mining hardware.}
Bitcoin mining chips are the SHA-256 circuits built in the largest numbers, but they are designed for energy per hash, not latency. Intel's Bonanza Mine~\cite{SKR+22} has 258 fully unrolled double-SHA-256 engines on a 7\,nm die. It runs at 355\,mV and 1.35--1.6\,GHz and computes 137 billion hashes per second at 55\,J per terahash. TICA~\cite{LHZ+22} pipelines the 64 rounds into 64 stages at 0.3\,V in 28\,nm. Both lower the supply voltage close to the threshold, where gates are several times slower, and recover throughput with many engines. Neither reports the latency of one hash. If each cycle of Bonanza Mine advanced one round, one compression would take 40--47\,ns, ten times our 3.5\,ns. A chain attacker would build the opposite: one engine at high voltage, with no registers between rounds.

\paragraph{The precedent of verifiable delay functions.}
Work on verifiable delay functions faced the same question and also answered it with a stated factor: Drake proposed $A_{\max} = 10$ as a conservative parameter for RSA-based VDFs, hoping that public ASICs would bring it to about 5~\cite{Drake18}. Here $A = 10$ means $\ell = 1.3$\,ns, about 4.5 FO4 per round, outside the model but not ruled out. \Cref{sec:open} lists a careful estimate as the most important open question.
